\chapter{Background}
\label{chap:background}

This chapter introduces the statistical framework for two-phase validation designs. It covers missingness mechanisms, estimation strategies for the resulting incomplete data, and sampling designs for selecting the phase-two subsample. The notation established here is used throughout the thesis unless stated otherwise.

\section{Electronic Health Records and the Need for Validation}
\label{sec:Background2.1}

EHR systems are widely used in U.S.\ clinical care. By 2021, $88.2\%$ of U.S.\ office-based physicians used an EHR system and $77.8\%$ used a certified system \citep{cdc_nehrs_2021_weighted}; $96\%$ of non-federal acute-care hospitals had adopted a certified EHR system \citep{onc2022report}. These systems contain records from large patient populations and repeated measurements collected during clinical care, which makes them useful for biomedical research.

EHR data were collected for clinical documentation and billing rather than research. Clinical notes use heterogeneous language across providers, clinics may apply different diagnostic guidelines and confirmation procedures, measurement protocols vary, and there is no designed recruitment process \citep{floyd2012use, beesley2020}. Without careful attention to variable definitions and observation windows, association analyses are vulnerable to information bias from error-prone variables. For a fixed measurement-error process, increasing the sample size reduces sampling error but does not remove bias in a regression coefficient \citep{carroll2006measurement}. The ``big data paradox'' of \citet{meng2018} concerns a related but distinct problem: in population averages, selection bias can dominate sampling error when inclusion is correlated with the variable being estimated.

The sources of information bias in EHR data are diverse. Diagnostic codes provide only a partial view of disease history; self-reported symptoms can be over-reported; secondary conditions and between-visit symptoms are frequently missed \citep{o2005}. Many studies model misclassification via constant sensitivity and specificity \citep{sinnott2014, lange2015, hubbard2015}, but in EHR data the error rate often depends on a patient's observation pattern, including the number of visits and the length of follow-up \citep{goldstein2016}. Approaches that allow sensitivity or specificity to vary with covariates without ground-truth labels have been proposed \citep{beesley2022statistical}, but they rely on identifying assumptions, such as conditional independence of misclassification errors across variables, that are difficult to verify and sensitive to model misspecification.

One well-established remedy is data validation, in which a trained auditor compares the derived EHR variables with the original source documents. Internal validation data can be used to correct analyses with error-prone continuous or discrete variables, and the approach has been applied in nutritional epidemiology and human immunodeficiency virus (HIV) cohort studies \citep{rosner1990correction, spiegelman1997fully, thurston2005comparison, lyles2011validation, shepherd2023multiwave}. In the seven-site HIV study of \citet{duda2012measuring}, $14\%$ of the fields examined in the initial audits contained major errors; the site-specific rates ranged from $2\%$ to $34\%$. Validation cannot recover information that was never recorded, and validating a full cohort is usually too costly. A probability subsample can instead be validated and analysed together with the unvalidated records \citep{huang2018pie, giganti2020accounting}.

In survey statistics, this subsample strategy is a two-phase design. Phase one records $(\tilde{X},Y,Z,A)$ for all $N$ individuals. Phase two obtains validated values of $X$ for $n$ individuals selected by probability sampling.

Let $Y = (Y_1, \ldots, Y_q)$ denote the outcome variables, $\tilde{X} = (\tilde{X}_1, \ldots, \tilde{X}_p)$ the error-prone unvalidated variables extracted from the EHR, $Z = (Z_1, \ldots, Z_k)$ the covariates that are error-free, and $A = (A_1, \ldots, A_l)$ any additional variables available at phase one. We assume that $Y$, $\tilde{X}$, $Z$, and $A$ are observed for every individual, while the validated variables $X = (X_1, \ldots, X_p)$, corresponding element-wise to $\tilde{X}$, are observed only in the phase-two subsample.

The key design quantity is the selection indicator. Let $R_i \in \{0,1\}$ indicate selection into phase two, with $\sum_{i=1}^{N} R_i = n$. If the subsample is drawn by simple random sampling, then $X$ is missing completely at random and the inclusion probability is $\pi_i = n/N$ for all $i$. If selection depends only on phase-one data, then $P(R \mid Y, \tilde{X}, Z, A, X) = P(R \mid Y, \tilde{X}, Z, A)$, so $X$ is missing at random \citep{rubin1976inference}, and the inclusion probability for individual $i$ is $\pi_i = E(R_i \mid Y_i, \tilde{X}_i, Z_i, A_i)$, with $\pi_i > 0$ required for all $i$. In some study designs, both the outcome $Y$ and the covariates $X$ are simultaneously error-prone, which still fits within the same framework \citep{giganti2020accounting}. Let $P(Y \mid X, Z;\, \beta)$ denote the outcome model with parameter vector $\beta$. We write $\beta_0$ for the true superpopulation value, where a superpopulation model is assumed, $\tilde{\beta}$ for the census estimator that would be obtained from the full population, and $\hat{\beta}$ for the estimator computed from the phase-two sample; when none is posited, the finite-population parameter $\tilde{\beta}$ is the estimand.

One reason for separating $Z$ and $A$ is that variables in $A$ may lie on the causal pathway from $X$ to $Y$. If a causal total effect is the target, conditioning on such a variable can block part of that pathway. The resulting coefficient has a direct-effect interpretation only under additional causal assumptions. Variables in $A$ can still inform the sampling design, the raking auxiliaries, the imputation model, or sensitivity analyses without being included in the outcome model. The surrogate $\tilde{X}$ is itself a variable of this kind: it informs the imputation model and the sampling design but does not enter the outcome model.


\section{Missingness Mechanisms}
\label{sec:missing_mech}

The standard classification of \citet{rubin1976inference} distinguishes three missingness mechanisms: missing completely at random (MCAR), missing at random (MAR), and missing not at random (MNAR). Under MCAR the probability of missingness is unrelated to both observed and unobserved values; under MAR it may depend on observed data but not on unobserved values conditional on what is observed; under MNAR it depends on unobserved values even after conditioning. Complete-case analysis is unbiased under MCAR but not in general. MAR together with distinct parameters for the data and missingness models gives the weakest general conditions under which likelihood-based or Bayesian inference can ignore the missingness process \citep{rubin1976inference, little2019statistical}. MAR is commonly used as a working assumption, with sensitivity analyses when MNAR is plausible.

The two-phase validation setting differs from a general missing-data problem because validation is assigned by design. Under the probability designs, each individual has a known positive inclusion probability $\pi_i = E(R_i \mid Y_i, \tilde{X}_i, Z_i, A_i)$ \citep{neyman1938contribution}. Simple random subsampling is MCAR, while probability sampling based on observed phase-one variables is MAR. The deterministic tail designs used later also select on observed phase-one information, but they are not probability samples: after the phase-one scores are fixed, individuals outside the selected region have no positive inclusion probability and no design weight. Their analysis therefore depends on a model that extrapolates the validated conditional relationship beyond the selected region. MAR and positive sampling support are separate conditions. If the phase-two sample also experiences nonresponse, the known design no longer accounts for all missingness, and additional response weighting or imputation is required.


\section{Two-Phase Design}
\label{sec:two_phase}

A two-phase design limits chart review to a subsample and combines the validated and unvalidated records in the analysis. For example, \citet{shepherd2023multiwave} validated 996 of $10{,}335$ mother-child EHR dyads. Error-prone outcomes such as obesity and asthma, together with covariates including sex, body mass index (BMI), maternal age, and maternal race, were available for all $10{,}335$ records; the validated versions were available only for the selected subsample.

\subsection{Likelihood Inference}
\label{sec:likelihood}

Following the notation of Section~\ref{sec:Background2.1}, the objective is to estimate the outcome-model parameters $\beta$. Let $P(X \mid \tilde{X}, Z, A;\, \alpha)$ denote the imputation model, with true parameter value $\alpha_0$ and sample-based estimator $\hat{\alpha}$. If $X$ were observed for all individuals, the complete-data likelihood would factor as
\begin{equation}
    L = \prod_{i=1}^{N}
    P(Y_i \mid X_i, \tilde{X}_i, Z_i, A_i)\;
    P(X_i \mid \tilde{X}_i, Z_i, A_i).
    \label{eq:complete_lik}
\end{equation}
Because $X_i$ is observed only when $R_i = 1$, the observed-data likelihood, conditional on the realised selection pattern, takes the form
\begin{equation}
    L = \prod_{R_i=1}
    P(Y_i \mid X_i, \tilde{X}_i, Z_i, A_i)\;
    P(X_i \mid \tilde{X}_i, Z_i, A_i)
    \;\prod_{R_i=0}
    P(Y_i \mid \tilde{X}_i, Z_i, A_i),
    \label{eq:obs_lik}
\end{equation}
where the phase-one contribution marginalises over the unobserved $X$:
\begin{equation}
    P(Y_i \mid \tilde{X}_i, Z_i, A_i)
    = \int P(Y_i \mid X, \tilde{X}_i, Z_i, A_i)\;
    P(X \mid \tilde{X}_i, Z_i, A_i)\, dX.
    \label{eq:marginal_Y}
\end{equation}
Strictly, \eqref{eq:obs_lik} omits the selection factor $\prod_i P(R_i \mid Y_i, \tilde{X}_i, Z_i, A_i)$, so it is proportional to, rather than equal to, the full likelihood. Under the ignorable design of Section~\ref{sec:missing_mech}, that factor is known and free of $(\beta, \alpha)$, so it drops out of their score. Inference for $\beta$ requires joint maximisation over $(\beta, \alpha)$, including numerical integration for the phase-one records. Consistency depends on correct specification of the outcome component $P(Y \mid X, \tilde{X}, Z, A)$ and the measurement component $P(X \mid \tilde{X}, Z, A)$ used in the likelihood.

In the measurement-error literature, \citet{prentice1982covariate} and \citet{breslow1997maximum} impose a pure surrogate assumption: $P(Y \mid X, \tilde{X}, Z, A) = P(Y \mid X, Z)$, which removes direct dependence on $(\tilde{X}, A)$ given $(X, Z)$. Under this assumption the observed-data likelihood simplifies to
\begin{equation}
    L = \prod_{R_i=1}
    P(Y_i \mid X_i, Z_i)\;
    P(X_i \mid \tilde{X}_i, Z_i, A_i)
    \;\prod_{R_i=0}
    \int P(Y_i \mid X, Z_i)\;
    P(X \mid \tilde{X}_i, Z_i, A_i)\, dX.
    \label{eq:obs_lik_pure}
\end{equation}
Design-based estimators do not require the pure surrogate assumption (Section~\ref{sec:DBE}). A likelihood analysis may instead allow direct dependence of $Y$ on $(\tilde{X},A)$ after conditioning on $(X,Z)$, but that dependence must be represented in the outcome model. A more flexible model for $P(X \mid \tilde{X}, Z, A)$ does not by itself correct a misspecified outcome factorisation \citep{prentice1982covariate, breslow1997maximum, breslow1999design}.

Phase-two selection may depend on the observed phase-one data, but its probability $\pi_i$ is fixed by the sampling design in the setting considered here. It contributes a known factor to the joint likelihood rather than a third fitted model. If selected records can also fail to respond, the response mechanism introduces an additional nuisance model, as described in Section~\ref{sec:missing_mech}; under the assumed absence of nonresponse, the unknown likelihood components are the outcome model indexed by $\beta$ and the measurement model indexed by $\alpha$.


\section{Design-Based Estimators}
\label{sec:DBE}

Design-based estimators specify only the sampling mechanism through the inclusion probabilities $\pi_i$, leaving the joint distribution of $(Y, X, Z)$ unrestricted. By reweighting the phase-two observations according to $\pi_i$, these estimators target full-cohort quantities without modelling $\tilde{X}$ directly. Consistency requires correct specification of $\pi_i$ and the positivity condition $\pi_i > 0$ for all $i$; unbiased variance estimation additionally requires the second-order (joint) inclusion probabilities $\pi_{ij} > 0$ for all pairs $(i,j)$. Asymptotic consistency may further require these probabilities to be bounded away from zero, $\pi_i \ge \pi_{\min} > 0$, so that they do not approach zero as $n$ increases.

\subsection{Inverse Probability Weighting}
\label{sec:ipw}

The inverse probability weighted (IPW) estimator, originally proposed by \citet{horvitz1952generalization} for estimating population totals under arbitrary unequal-probability sampling, estimates $\beta$ by maximising the weighted log-likelihood
\begin{equation}
    \ell_{\mathrm{IPW}}(\beta)
    = \sum_{i=1}^{N} \frac{R_i}{\pi_i}\,
    \log f(Y_i \mid X_i, Z_i;\, \beta),
    \label{eq:ipw}
\end{equation}
so that only units with $R_i = 1$ contribute, each through the weight $1/\pi_i$. In the survey literature this estimator is often referred to as the pseudo-likelihood estimator \citep{pfeffermann1993role}. For a correctly specified complete-data estimating equation, the design-weighted IPW estimator is consistent when the correct inclusion probabilities are used, without requiring a model for $X \mid Z$. Its precision depends on the sampling design and the phase-two sample size. A high variance of the IPW estimator does not by itself imply inefficiency: under some sampling designs it is fully efficient. The efficiency it does forgo comes from using only the selected units in the score of~\eqref{eq:ipw} and discarding the phase-one information available on the unvalidated records, which the augmented estimators of the following sections recover. Generalised raking can reduce its variance when phase-one auxiliaries predict the complete-data score contributions.

\subsection{Generalised Raking}
\label{sec:raking}

Calibration of weights, or generalised raking, proposed by \citet{deville1992calibration}, improves upon IPW by adjusting the design weights using phase-one information \citep{breslow2009using}. The name \emph{calibration} reflects the defining requirement that the adjusted weights reproduce (calibrate to) known phase-one totals of auxiliary variables when applied to the phase-two sample. Throughout this thesis we use the term ``raking'' to avoid ambiguity with other uses of ``calibration'' in the measurement-error literature.

The raking estimator maximises the same pseudo-likelihood as IPW but with calibrated weights $w_i = g_i / \pi_i$:
\begin{equation}
    \ell_{\mathrm{raking}}(\beta)
    = \sum_{i=1}^{N} \frac{R_i\, g_i}{\pi_i}\,
    \log f(Y_i \mid X_i, Z_i;\, \beta),
    \label{eq:raking}
\end{equation}
where the rescaling factors $g_i$ are obtained by solving
\begin{equation}
    \min_{\{g_i\}} \sum_{i=1}^{N} R_i\,
    d\!\left(\frac{g_i}{\pi_i},\; \frac{1}{\pi_i}\right)
    \quad \text{subject to} \quad
    \sum_{i=1}^{N} R_i\, \frac{g_i}{\pi_i}\, H_i
    = \sum_{i=1}^{N} H_i,
    \label{eq:raking_opt}
\end{equation}
and $d(\cdot, \cdot)$ is a smooth distance function and $H$ denotes a vector of fully observed phase-one auxiliaries. The objective minimises the deviation of the calibrated weights from the base weights, while forcing the weighted phase-two totals of $H$ to match the known phase-one totals. Two common choices for the distance function are the quadratic form $d_1(a,b) = (a - b)^2 / (2b)$, which yields the generalised regression (GREG) estimator, and the entropy form $d_2(a,b) = a\log(a/b) + b - a$, which yields the multiplicative raking estimator \citep{lumley2011connections}. The quadratic distance can produce negative calibrated weights unless explicitly constrained, whereas the entropy distance produces positive weights by construction. This matters in practice: many regression and survey routines reject negative weights or handle them incorrectly, so the positive weights of $d_2$ can be used directly with such software. \citet{deville1992calibration} showed that the GREG estimator is the first-order approximation to generalised raking under any distance function, so that variance estimators developed for GREG apply broadly. In the present thesis we adopt $d_2$. The optimisation in~\eqref{eq:raking_opt} is solved by introducing Lagrange multipliers for the calibration constraints and applying Newton--Raphson iteration; the resulting rescaling factors $g_i$ take a generalised-linear form, with the link function determined by the choice of distance $d(\cdot,\cdot)$.

The efficiency gain from generalised raking depends on how well the auxiliary variables predict the complete-data estimating function. For a population total, the relevant residual is obtained by projecting the study variable on the calibration auxiliaries \citep{deville1992calibration}. For a regression coefficient, the corresponding quantity is its complete-data score or influence function, and calibrating on a close phase-one approximation can reduce variance substantially \citep{lumley2011connections}. A phase-one proxy that is strongly associated with a phase-two variable can help, but that association matters through its prediction of the estimating function rather than by itself.

If $\tilde{\beta}$ is regular and asymptotically linear, then, by definition,
\begin{equation}
    \sqrt{N}\,(\tilde{\beta} - \beta_0)
    = \frac{1}{\sqrt{N}} \sum_{i=1}^{N} \mathrm{h}_i(\beta_0)
    + o_p(1),
    \label{eq:al_rep}
\end{equation}
where $\mathrm{h}_i(\beta_0)$ is the influence-function contribution for unit $i$. The influence function can be derived from the estimating equation $\sum_{i=1}^{N} U_i(\beta) = 0$, where $U_i(\beta)$ is the score contribution from individual $i$. Defining the information matrix $\mathcal{I} = -\mathrm{E}[\partial U_i(\beta_0) / \partial \beta]$, the influence function is $\mathrm{h}_i(\beta_0) = \mathcal{I}^{-1} U_i(\beta_0)$ under standard regularity conditions. In practice, it is evaluated at an estimated coefficient and can be approximated via the \texttt{dfbeta} method \citep{therneau2000cox} or via the score function directly.

\citet{breslow2009improved} showed that generalised raking achieves its greatest efficiency when the raking auxiliaries are highly correlated with $\mathrm{h}_i(\beta)$, because the estimation error $\tilde{\beta}-\beta_0$ has a first-order representation as the mean of the influence-function contributions. The optimal auxiliary is
\begin{equation}
    H_i = \mathrm{E}\bigl[\mathrm{h}_i(\beta)
    \mid Y_i, \tilde{X}_i, Z_i, A_i\bigr],
    \label{eq:optimal_aux}
\end{equation}
and generalised raking with this choice attains the optimum within the corresponding augmented inverse probability weighted class \citep{lumley2011connections}. Under quadratic calibration, projecting the influence function onto the phase-one auxiliaries gives a raked score that is algebraically identical to the corresponding augmented inverse probability weighted (AIPW) equation. For the entropy distance used in this thesis, the two equations agree to first order rather than exactly. Appendix~\ref{app:calib-aipw} states the quadratic identity and the first-order qualification for other calibration distances.

In two-phase validation studies, $X$ and $\tilde{X}$ are often strongly correlated, so replacing $X$ with $\tilde{X}$ in $\mathrm{h}_i(\beta)$ gives a practical proxy for $H_i$. A model for $X \mid (Y, \tilde{X}, Z, A)$ fitted to the phase-two data can also be used to impute $\hat{X}$ and approximate $\mathrm{E}[\mathrm{h}_i(\beta) \mid Y, \tilde{X}, Z, A]$ \citep{kulich2004improving}; either single or multiple imputation can construct this auxiliary \citep{han2021combining}. For a fixed auxiliary, misspecifying that model changes the efficiency gained from calibration rather than the design consistency of the raked estimator. First-order robustness to estimating the auxiliary is not automatic: it requires the calibration estimating equation to remain unbiased for each fixed nuisance value, together with smoothness, orthogonality and a nuisance rate faster than $n^{-1/4}$. Under these conditions, estimating the auxiliary does not change the consistency or first-order variance of $\hat{\beta}$. Appendix~\ref{app:fourth-root} gives the general result and verifies its use for the calibration equation.

In this thesis, the generalised raking estimator is implemented using the \texttt{survey} package \citep{lumley2011complex} in \texttt{R} \citep{r2025}. The procedure is as follows. First, imputations $\hat{X}$ are obtained for all individuals by fitting weighted imputation models on the observed phase-two data. Second, the influence function $\mathrm{h}_i(\beta)$ is computed using the imputed values to fit the outcome model; for linear regression, the \texttt{dfbeta} approximation in \texttt{R} is used, while for other regression families the score function is evaluated directly. Third, the calibrated weights are obtained using the entropy distance $d_2$. Fourth, $\beta$ is estimated by maximising the raking pseudo-likelihood \eqref{eq:raking} with the calibrated weights.


\section{Model-Based Estimators}
\label{sec:model_based}

Model-based estimators use distributional assumptions about the validated variables, the outcome, or both. They can gain precision by combining the phase-one and phase-two records when those assumptions are adequate, while misspecification can introduce bias. This section reviews the principal model-based approaches used in two-phase validation studies.

\subsection{Plug-In Likelihood}
\label{sec:plugin}

Plug-in likelihood estimation addresses the two-phase problem by maximising a log-likelihood in which nuisance parameters are estimated empirically. For logistic regression under case-control sampling, \citet{breslow1988logistic} proposed a two-step estimator that first estimates the nuisance parameters and then plugs them into the log-likelihood to estimate $\beta$. \citet{schill1993logistic} proposed a related approach that estimates $\beta$ and the nuisance parameters simultaneously. Both methods assume the pure surrogate condition $P(Y \mid X, \tilde{X}, A) = P(Y \mid X)$. Using the National Wilms' Tumor Study data, \citet{breslow1999design} showed that the two approaches perform similarly and are more efficient than IPW when the pure surrogate assumption holds. \citet{rivera2020augmented} extended these methods to augmented plug-in likelihood estimators that exploit whole-cohort information.

\subsection{Semiparametric Maximum Likelihood}
\label{sec:spmle}

The semiparametric maximum likelihood estimator (SPMLE) estimates the outcome-model parameters by maximising the likelihood while estimating $P(X \mid Z)$ nonparametrically from the data. In practice, $P(X \mid Z)$ is represented by a discrete distribution supported on the distinct observed values of $X$. Under correct specification of the outcome model, the resulting estimator attains full semiparametric efficiency \citep{breslow1997maximum, scott1997fitting}. Under mild misspecification, the SPMLE can become biased and lose its efficiency advantage, so design-based methods may offer better robustness \citep{han2021combining}.

For outcome-dependent sampling at phase two with discrete $X$ and simple random sampling at phase one, \citet{scott1997fitting} showed that the SPMLE can be constructed by iterating conditional maximum likelihood with an updating offset. Dimension control is important: if $X$ is continuous, profiling introduces arrays of dimension $J - 1$, where $J$ is the number of distinct values of $X$. \citet{ScottWild2006-efficient-spmle} address this by collapsing replicated values. When the outcome $Y$ is discrete, a reformulation of the profile likelihood reduces the profiling dimension to $K - 1$, where $K$ is the number of outcome categories.

\citet{tao2017efficient} generalised the approach to accommodate both continuous and discrete outcomes by approximating $P(X \mid Z)$ via sieve approximations \citep{Grenander1981-AbstractInference} built from B-spline bases \citep{Schumaker1981-SplineFunctions} and maximising the resulting sieve likelihood via an EM algorithm. The number of sieve terms increases with the phase-one sample size $N$, so computational cost grows accordingly, and the curse of dimensionality constrains the number of continuous features that can be included in $Z$.

\subsection{Regression Calibration}
\label{sec:regcal}

Regression calibration \citep{prentice1982covariate} is a standard approach for handling measurement error in covariates. It replaces the unobserved $X$ in the outcome model with its conditional expectation given the observed data, yielding a working model $P(Y \mid \mathrm{E}(X \mid \tilde{X}), Z)$. Under the classical additive error model
\begin{equation}
    \tilde{X} = X + S, \qquad
    \mathrm{E}(S) = 0, \quad
    \mathrm{Var}(S) = \sigma_S^2,
    \label{eq:me}
\end{equation}
if $X$ and $S$ are independent normal variables and $\mathrm{E}(X)=0$, then $\mathrm{E}(X \mid \tilde{X}) = \lambda\, \tilde{X}$, where $\lambda = \mathrm{Var}(X) / \mathrm{Var}(\tilde{X})$ is the reliability ratio. Without normality, the same expression is the best linear predictor when $X$ and $S$ are uncorrelated, but it need not equal the conditional expectation. In practice, the calibration equation is estimated from the validation data. Standard errors that treat its estimated parameters as fixed omit their estimation uncertainty; a joint sandwich calculation or resampling procedure can account for it \citep{huber1967behavior, white1980heteroskedasticity}.

For a linear outcome model, regression calibration consistently estimates $\beta$ when the error is nondifferential and the calibration equation $\mathrm{E}(X \mid \tilde{X},Z)$ is correctly specified. For non-linear outcome models such as logistic regression, replacing $X$ by its conditional expectation is generally approximate and can induce bias, though the bias can be small when the measurement-error variance is moderate relative to the signal variance \citep{carroll2006measurement}.

\subsection{Multiple Imputation}
\label{sec:mi}

Multiple imputation \citep{Rubin1987} replaces the missing values by draws from the posterior predictive distribution, producing $m$ completed datasets that each support a standard complete-data analysis. In the two-phase validation setting, write $X = (X_{R=0}, X_{R=1})$, where $X_{R=1}$ are the phase-two measurements and $X_{R=0}$ the unobserved values for unselected individuals. Let $\mathcal{D}_{\mathrm{obs}}=(Y,\tilde{X},Z,A,R,X_{R=1})$ denote all observed data. The posterior for $\beta$ can then be expressed as
\begin{equation}
    P(\beta \mid \mathcal{D}_{\mathrm{obs}})
    = \int P(\beta \mid \mathcal{D}_{\mathrm{obs}}, X_{R=0})\;
    P(X_{R=0} \mid \mathcal{D}_{\mathrm{obs}})\, dX_{R=0}.
    \label{eq:mi_posterior}
\end{equation}
Operationally, $X_{R=0}$ is replaced with $m$ draws from an approximation to the posterior predictive distribution. Fitting the analysis model to each completed dataset gives estimates $(\hat{\beta}^{(1)}, \ldots, \hat{\beta}^{(m)})$ and variance estimates $(\hat{V}^{(1)}, \ldots, \hat{V}^{(m)})$. The combined coefficient estimate is
\begin{equation}
    \bar{\beta} = \frac{1}{m} \sum_{j=1}^{m} \hat{\beta}^{(j)},
    \label{eq:mi_point}
\end{equation}
and the total variance is
\begin{equation}
    \hat{V}_{\mathrm{total}}
    = \hat{V}_{\mathrm{within}}
    + \frac{m+1}{m}\, \hat{V}_{\mathrm{between}},
    \label{eq:mi_var}
\end{equation}
where $\hat{V}_{\mathrm{within}} = m^{-1} \sum_{j=1}^{m} \hat{V}^{(j)}$ reflects the average complete-data variance and $\hat{V}_{\mathrm{between}} = (m-1)^{-1} \sum_{j=1}^{m} (\hat{\beta}^{(j)} - \bar{\beta})(\hat{\beta}^{(j)} - \bar{\beta})^{\top}$ measures variation between the completed-data estimates. The factor $(m+1)/m$ is the standard finite-$m$ correction from Rubin's rules \citep{Rubin1987}.

Unlike single imputation and regression calibration, multiple imputation uses within-imputation and between-imputation components to represent uncertainty from the unobserved data. Rubin's variance gives valid frequentist inference when the imputation is proper for the complete-data analysis and that analysis is itself valid; congeniality with a correctly specified embedding model provides sufficient conditions \citep{meng1994multiple}. Under uncongeniality or model misspecification, Rubin's variance can be biased upward or downward. The conservative intervals studied by \citet{fay1992inferences} are one case, not a general direction of bias.

The imputation model should represent the transformations, interactions and design variables needed by the analysis unless stronger restrictions are justified \citep{meng1994multiple}. The number of auxiliary variables must still be judged against the observed sample size, because an overly large imputation model can reduce precision and bias regression coefficients in small samples \citep{hardt2012auxiliary}. Bootstrap validity depends on the order of the two procedures: imputation followed by bootstrap is generally not valid under uncongeniality or misspecification, whereas specified bootstrap-then-imputation procedures can be valid \citep{bartlett2020bootstrap}. The variance estimator of \citet{robins2000inference} avoids a general congeniality requirement but needs estimating equations for the particular imputation and analysis procedures. It therefore requires a method-specific derivation for chained equations or machine-learning imputations rather than transferring automatically \citep{hughes2016comparison}.

Two broad strategies for multivariate imputation are joint modelling \citep{carpenter2023multiple} and fully conditional specification (FCS) \citep{vanBuuren2006}. Joint modelling specifies one multivariate distribution and draws missing values from its posterior predictive distribution; its performance depends on that distributional specification. FCS cycles through separately specified conditional models. It can be more robust than a misspecified restricted joint model in some settings, but the conditional models need not correspond to any joint distribution \citep{seaman2018relative, white2011multiple}. Both strategies are reviewed in Chapter~\ref{chap:ImpReview}.

In the two-phase applications considered here, imputations for $X_{R=0}$ condition on the available phase-one information $(Y, \tilde{X}, Z, A)$. When this set contains every variable used for selection, the designed missingness is MAR for $X$. Conditioning on those variables is not itself a design-based correction: an imputation procedure does not automatically account for unequal inclusion probabilities, and the design variables or weights must be used when they carry information not otherwise represented. The same outcome model is fitted to the $m$ completed datasets and the estimates are combined by Rubin's rules. Multiple imputation can also construct phase-one proxies for the influence functions used in raking; under the conditions of Section~\ref{sec:raking}, their accuracy governs the efficiency gained from calibration.


\section{Sampling Designs}
\label{sec:sampling}

Simple random sampling is widely used because it is straightforward to implement and analyse. However, fully random selection ignores the information contained in the phase-one auxiliary variables, which may be highly correlated with the unobserved truth. Returning to the error model \eqref{eq:me}, the unvalidated variable $\tilde{X}$ carries information about $X$ that can be exploited by the sampling design to reduce the variance of $\hat{\beta}$. An optimal sampling mechanism selects the validation subsample so as to minimise $\mathrm{Var}(\hat{\beta})$ for a given budget $n$ \citep{amorim2021two}. Following the optimal-design framework of \citet{amorim2021two}, the designs reviewed below are grouped by the class of estimator they are tailored to: \emph{design-based} sampling, aimed at the design-based (IPW and raking) estimators, and \emph{model-based} sampling, aimed at the model-based (likelihood) estimators.

\subsection{Design-Based Sampling}
\label{sec:design_sampling}

Design-based sampling strategies must satisfy the positivity condition: $\pi_i > 0$ for all individuals. Write $k$ for the stratum index, $N_k$ for the stratum population size, $n_k$ for the number selected from stratum $k$, and $P_k=N_k/N$ for the stratum's population share.

\paragraph{Stratified sampling.}
Stratified sampling can reduce the variance of $\hat{\beta}$ when the strata are homogeneous in the quantities that drive its influence function and the sample is allocated appropriately. Under the framework of Section~\ref{sec:DBE}, the variance decomposes into within-stratum contributions. Under proportional allocation, the standard variance comparison depends on the average within-stratum variance rather than the total variance, while Neyman allocation also uses differences among the within-stratum standard deviations. An arbitrary allocation need not improve on \texttt{SRS}; the gain depends on the allocation and on how well the stratification variable separates the relevant influence contributions \citep{cochran1977, sarndal1992model, lumley2011complex}.

\paragraph{Neyman allocation.}
Given the stratum structure, the within-stratum sample sizes $n_k$ can be optimised to minimise the variance of $\hat{\beta}$. For a scalar estimand, the allocation-dependent part of the stratified IPW variance is $\sum_{k=1}^{K} P_k^2 s_k^2 / n_k$, where $s_k$ is the within-stratum standard deviation of the influence function $\mathrm{h}_i(\beta)$. The optimisation problem is
\begin{equation}
    \min_{\{n_k\}} \; \sum_{k=1}^{K} \frac{P_k^2\, s_k^2}{n_k} + \lambda \biggl( \sum_{k=1}^{K} n_k - n \biggr),
    \label{eq:neyman_lagrange}
\end{equation}
where $\lambda$ is a Lagrange multiplier enforcing the budget constraint $\sum_{k=1}^{K} n_k = n$. Minimising with respect to $n_k$ yields the Neyman allocation
\begin{equation}
    \frac{n_k}{n} = \frac{N_k\, s_k}{\sum_{k'=1}^{K} N_{k'}\, s_{k'}},
    \label{eq:neyman_alloc}
\end{equation}
\citep{neyman1938contribution}. Strata with larger populations or greater variability in the influence function receive proportionally more of the phase-two budget. The continuous Neyman allocation is converted to integer stratum sample sizes using the exact optimal algorithm of \citet{wright2017exact}, as described in Appendix~\ref{app:sim1_sampling}.

\paragraph{Optimal strata boundaries.}
Once the stratification variable and the number of strata $K$ are chosen, the boundaries can be selected to reduce the variance of the scalar IPW target. Under Neyman allocation, the variance is proportional to $\left(\sum_{k=1}^{K} N_k s_k\right)^2$; the construction used by \citet{amorim2021two} chooses boundaries for which $N_k s_k$ is approximately constant, leading to approximately equal stratum allocations. In their linear measurement-error simulations, moving from one to three influence-function strata produced the largest gain, with smaller additional gains from further strata. The three-stratum choice is therefore adopted as the default for the stratified model-based designs used in Simulation Design~II and the VCCC case study; Simulation Design~III instead uses the cross-classified strata defined for its survival setting and does not follow this rule.

\subsection{Model-Based Sampling}
\label{sec:model_sampling}

\paragraph{Outcome-dependent sampling.}
Outcome-dependent sampling (\texttt{ODS}) selects phase-two units on the basis of the outcome $Y$. For discrete $Y$, the prototypical example is case-control sampling \citep{breslow1988logistic}: with a fixed validation budget $n$, \texttt{SRS} produces approximately $n \cdot p_{\mathrm{prev}}$ cases when the prevalence $p_{\mathrm{prev}}$ is small, whereas a balanced case-control design allocates roughly $n/2$ cases and $n/2$ controls. Oversampling cases therefore provides many more case-control comparisons than \texttt{SRS} when the outcome is rare, which can improve precision for odds-ratio parameters \citep{prentice1979logistic}. The efficient case fraction depends on the covariate distribution and the target parameters; equal numbers of cases and controls are useful in many settings but are not universally optimal \citep{cain1988twoStage}.

For continuous $Y$, sampling from the tails of the outcome distribution places validation effort in regions that exert the greatest influence on the regression parameters. \citet{zhou2002semiparametric} showed that conditioning on extreme values of $Y$ reduces the asymptotic variance of $\hat{\beta}$ under a bivariate normal model, and semiparametric estimators can attain the efficiency bound under two-stage \texttt{ODS}. In practice, one may discretise $Y$ into $K$ strata and oversample the tails, or rank $Y$ and select the $n/2$ individuals with the largest and $n/2$ with the smallest values \citep{zhou2002semiparametric, song2009twostage}.

\paragraph{Residual sampling.}
Residual sampling (\texttt{RS}) selects phase-two units whose outcomes are poorly predicted by the phase-one covariates. A phase-one regression of $Y$ on $(Z, \tilde{X})$ is fitted, residuals $\hat{e}_i$ are computed for all $N$ individuals, and the $n/2$ individuals with the most extreme positive residuals and $n/2$ with the most extreme negative residuals are selected deterministically. The rationale is that individuals whose outcomes are poorly predicted by the phase-one model carry the most information about the unobserved reference value $X$ \citep{breslow2009improved}. \citet{amorim2021two} showed that residual sampling performs similarly to the locally optimal weighted-residual design when $\beta$ is close to zero and $\operatorname{Var}(X\mid\tilde{X},Z)$ is constant across $Z$. The residuals use only phase-one variables, so selection can be completed before validation.

\paragraph{Weighted residual sampling.}
\citet{tao2019optimal} derived an optimal model-based selection criterion under the local alternative $\beta = o(1)$. For a normal outcome, the criterion is $(Y_i - \mu(Z_i))^2 \operatorname{Var}(X_i \mid \tilde{X}_i, Z_i)$, where $\mu(Z_i)$ is the linear predictor from a regression of $Y$ on $Z$ alone. Weighted residual sampling (\texttt{WRS}) ranks the signed quantity $(Y_i-\mu(Z_i))\{\operatorname{Var}(X_i\mid\tilde{X}_i,Z_i)\}^{1/2}$ and selects observations from its two tails. If the conditional variance is constant, this ranking is the same as ordinary residual sampling \citep{amorim2021two}. The variance term must be obtained from external information or a pilot subsample, which reduces the remaining targeted-sample budget.

This design is optimal when $\beta$ is close to the null hypothesis, but its performance can degrade as $\beta$ deviates from zero. In the simulations of \citet{amorim2021two}, weighted residual sampling and ordinary residual sampling produced nearly identical efficiency gains when $\operatorname{Var}(X\mid\tilde{X},Z)$ was constant across $Z$. The need for a pilot sample and the sensitivity to the local-alternative assumption are the main practical limitations of this design.

\paragraph{Score function sampling.}
\citet{amorim2021two} proposed an ad hoc design that avoids the pilot-sample requirement of weighted residual sampling. At $\beta=0$ in the normal linear model, the complete-data score for $\beta$ is proportional to $(Y_i - \mu(Z_i)) X_i$, and its conditional second moment gives the weighted-residual criterion. Score function sampling (\texttt{SFS}) replaces the unobserved $X_i$ in that score with the phase-one surrogate $\tilde{X}_i$, then selects the $n/2$ largest and $n/2$ smallest values of $(Y_i - \hat{\mu}(Z_i)) \tilde{X}_i$. Its usefulness depends on how well $\tilde{X}_i$ represents $X_i$ \citep{amorim2021two}.

In the simulations of \citet{amorim2021two}, \texttt{SFS} outperformed residual and weighted residual sampling when $\beta$ was away from zero and $\tilde{X}$ was at least a fair surrogate for $X$. Near the null, weighted residual sampling was more efficient. When no pilot information is available, the choice between residual sampling and \texttt{SFS} depends on prior information: residual sampling is favoured near the null or with a poor surrogate, whereas \texttt{SFS} is favoured away from the null with a fair surrogate.

\paragraph{Influence function sampling.}
For design-based estimators, the variance of $\hat{\beta}$ is asymptotically proportional to the variance of the influence function contributions $\mathrm{h}_i(\beta)$ (Section~\ref{sec:raking}). A natural stratification variable is therefore the influence function itself, and the optimal within-stratum allocation follows from the Neyman criterion~\eqref{eq:neyman_alloc}. Because $X$ is unobserved at phase one, the true influence function cannot be computed, and one instead substitutes $\tilde{X}$ for $X$ in the score and information matrix to obtain an approximate influence function $\tilde{\mathrm{h}}_i(\beta)$. Strata are then formed by partitioning $\tilde{\mathrm{h}}_i(\beta)$ into $K$ groups, and the phase-two allocation follows $n_k \propto N_k s_k$, where $s_k$ is the within-stratum standard deviation of $\tilde{\mathrm{h}}_i(\beta)$. This is the influence function sampling (\texttt{IFS}) design of \citet{amorim2021two}.

The quality of \texttt{Neyman-IF} depends on how well the surrogate-based influence function $\tilde{\mathrm{h}}_i(\beta)$ approximates the true influence function $\mathrm{h}_i(\beta)$. \citet{sarndal1992model} suggested a rule of thumb that the correlation between the stratification variable and the target quantity should be at least $0.90$ for efficient stratification. In the strongest-surrogate setting of \citet{amorim2021two}, influence-function sampling gave the smallest IPW variance among the designs compared. Its advantage diminished as the surrogate weakened, and other stratification variables performed better in the weakest-surrogate setting.

Raking already uses the phase-one influence function as an auxiliary for calibrating the design weights, so stratifying on the same quantity can be partly redundant. In the simulations of \citet{amorim2021two}, raking under IFS was sometimes less efficient than raking under residual-based stratification even when the surrogate was strongly correlated with the truth. The result indicates that efficiency also depends on what the stratification variable contributes beyond the calibration auxiliary.

\subsection{Extension to Survival Outcomes}
\label{sec:sampling_survival}

The designs described above were developed for the linear regression setting. \citet{han2021two} extended the mean score method of \citet{Reilly1995} to discrete-time survival models and derived an analogous Neyman-type optimal allocation for the proportional hazards setting. The discrete-time complementary log-log model has the same regression coefficient $\beta$ as the continuous-time Cox model, so a design derived for the discretised outcome can also improve efficiency for the continuous-time analysis. In their simulation study, \citet{han2021two} found that the adaptive mean score design, which first draws a balanced pilot sample to estimate the nuisance parameters and then allocates the remaining budget according to the estimated optimal proportions, produced consistent efficiency gains over \texttt{SRS} and balanced sampling when the phase-two sample size was at least moderately large. When the phase-two sample was very small, \texttt{SRS} was preferred because the pilot was too imprecise to estimate the nuisance parameters reliably.

When the survival outcome is itself error-prone, the surrogate outcome must replace the true outcome in both the stratum definitions and the phase-one model, following the dependent-error framework of \citet{giganti2020accounting}. This joint covariate-and-outcome error structure imposes additional demands on the sampling design, because the stratification variable is itself measured with error and the resulting strata may not align well with the true influence function.

\subsection{Summary}

\texttt{SRS} is the baseline. For fixed strata and a scalar IPW target, Neyman allocation minimises the asymptotic variance over allocations of a fixed phase-two sample. Outcome-dependent, residual and weighted residual sampling instead target model-based estimators, while influence-function sampling targets an IPW variance contribution. The best design depends on the estimator, and a design chosen for IPW need not be optimal for raking or likelihood analysis \citep{amorim2021two}. The simulations therefore evaluate both design-based and imputation-based estimators under several distinct sampling designs.
